\chapter{Settings, preferences and storage}\label{ch:storage}

\section{The Settings page}
Besides the device form (Chapter~\ref{ch:connect}) the \ui{Settings} page has an \ui{App} card with two items.
\begin{description}[leftmargin=3.2cm,style=nextline]
\item[Time shown as] \ui{UTC} or \ui{Local}. It changes the time zone of the fields and axes on the Live, Data and Helicorder pages; stored data
  and exports are always UTC. The same toggle is on the Data page.
\item[Data folder] the folder where everything is kept (Section~\ref{sec:folder}).
\end{description}

\section{What the application remembers}
The application remembers, between sessions, the Shake you were using, your choice of UTC or local time, and your last choices on
the data pages: the station, \ui{counts} or \ui{m/s}, whether the band-pass was on, and its two frequencies. These are saved the
moment you change them. Passwords are not saved in any file: they go to your operating system's key store.

\section{Inside the data folder}\label{sec:folder}
\begin{center}
\begin{tabular}{@{}L{0.22\linewidth}L{0.68\linewidth}@{}}
\toprule
Item & Contents\\ \midrule
\file{sds/} & your data: one miniSEED file per channel per day, in the SDS layout (Appendix~\ref{app:formats})\\
\file{response/} & instrument responses (StationXML) downloaded for the \ui{m/s} switch\\
\file{config.toml} & the list of Shakes (name, address, ports, user; no passwords)\\
\file{settings.toml} & the preferences described above\\
\file{known\_hosts} & the SSH identities of your Shakes\\
\file{logs/} & log files\\
\bottomrule
\end{tabular}
\end{center}

\paragraph{Back up.} Copy the whole \file{sds/} folder (and, if you like, \file{config.toml}). It contains ordinary files that other
seismology software can read directly.

\paragraph{Move to another disk.} Close the application, copy the folder, and start the application with \code{RS3D\_DATA\_DIR} set to
the new place (Section~\ref{ch:install}).

\paragraph{Reset the stored Shake identity.} If a Shake's SD card is replaced, its SSH identity changes and the \ui{SSH} check reports it.
Delete the line of that address from \file{known\_hosts} and refresh the Status page.

\paragraph{Free space.} A day of three channels occupies about 85~MB, so a year is about 30~GB.

\begin{rs3dwarn}[The Shake keeps only a week]
The Shake deletes days older than its retention (7 days by default). Download new days regularly (the Status page warns when days
are about to be deleted); your laptop's copy may then be the only one.
\end{rs3dwarn}
